%Chapter 1

\renewcommand{\thechapter}{1}

\chapter{Introduction}\label{chap:intro}

Software is written with intent: a programmer wants a computer to do something.
Software engineering is the process of refining and reformulating that intent into a
form the computer can execute. Software testing is the bridge between intent and the
produced program. It is the process of trying inputs, measuring behavior, and
understanding what the program does separately from how it is implemented.

%% Source for arguments and vocabulary:
%% https://alperenkeles.com/posts/vocab-for-testing/

Different methods of software testing ask the programmer to express intent in
different forms. In the simplest scenario, programmers \emph{tinker} with the
program at hand. They interact with the produced artifact, pressing buttons in a user
interface, exploring different options in a command-line interface, testing a variety
of inputs for a function through a REPL, and observing the results. These interactions
are then \emph{codified}, written as concrete testing instructions to be run alongside
program development as part of an automated test suite. The simplest form of
codification is \emph{example-based testing}, in which the specification relates
concrete user-written input-output pairs. A sorting function is tested through several assertions
such as \code{sort([]) == []}, \code{sort([1, 0]) == [0, 1]}, \code{sort(3,2,-1) ==
[-1, 2, 3]}.

Another form of specification, the one that will be the focus of this thesis, is
\emph{properties}. A property is a universally quantified predicate function: a fact
about the program under test that must hold for any possible input. Since properties
are universal, they let users specify whole classes of behavior instead of individual
examples. Once a property is written as a specification of a particular program
behavior, there are different strategies for producing inputs: systematic exhaustive
search, random sampling, and perhaps even user-supplied inputs. The combination of
these methods is typically gathered under the umbrella of \emph{property-based
    testing}.

Property-based testing (PBT) was pioneered by Koen Claessen and John Hughes in their
ICFP 2000 paper ``QuickCheck: a lightweight tool for random testing of Haskell
programs''~\cite{ClaessenH00}. The early history of random software testing is
typically traced to the CS736 class at the University of Wisconsin--Madison, where Bart
Miller coined the term ``fuzz'' testing by fuzzing UNIX utilities with random inputs to
trigger faults and find bugs~\cite{Miller1990UnixFuzz}. While fuzz testing has evolved
toward increasingly feedback-guided, target-aware, and domain-specific input generation
techniques~\cite{fuzzingbook2024,aflplusplus,DBLP:journals/tse/BohmePR19,FairFuzz,DirectedGreybox,FuzzFactory},
property-based testing has chosen specifications as its main focus. The original paper
introduces both a language for expressing properties and a library implementation for
writing random data generators. This design has been replicated in virtually all
programming languages in the last 26 years.
Table~\ref{tab:pbt-frameworks} gives a glimpse of QuickCheck's reach in the random
testing world by presenting the community-maintained
\code{jmid/pbt-frameworks} survey~\cite{jmid-pbt-frameworks}.

The ubiquity of the approach stems from its simplicity. The user starts with a simple
universal predicate denoting the correctness of their program:

\begin{code}[language=Haskell,emph={sort_correct,is_sorted,sort,is_permutation},emph={[2]l}]
sort_correct :: Ord a => [a] -> Bool
sort_correct l = is_sorted (sort l) && is_permutation l (sort l)
\end{code}

The \code{sort_correct} predicate denotes the correctness of a sorting function, the
elements of the list must not change, and the output must obey an ordering relation.
The user then writes a random data generator for the input type. Since
\code{sort_correct} is defined over arbitrary lists of orderable elements, we need a
\code{List} generator. For this example, we can instantiate the orderable element type with
\code{Int}s. QuickCheck provides a library of combinators for writing
generators such as \code{choose} and \code{listOf}, and the \code{Arbitrary} typeclass
for denoting generators of certain types.

\begin{code}[language=Haskell,emph={arbitrary,choose,listOf},emph={[2]a}]
instance Arbitrary Int where
    arbitrary = choose (0, 100)

-- listOf is a combinator that takes a generator for a
-- type and produces a generator for lists of that type.
instance Arbitrary a => Arbitrary [a] where
    arbitrary = listOf arbitrary
\end{code}

Once the universally quantified predicate and the corresponding random data generators
are in place, we can use QuickCheck's \code{forall} combinator for combining them into
properties to be tested.

\begin{code}[language=Haskell,emph={prop_sort_correct,forAll,arbitrary,sort_correct,quickCheck},emph={[2]xs}]
prop_sort_correct :: Property
prop_sort_correct = forAll arbitrary $ \xs -> sort_correct xs

quickCheck prop_sort_correct
{-
Stdout:
+++ OK, passed 100 tests.
-}
\end{code}

%% Source:
%% Programmable Property-Based Testing
%% https://arxiv.org/pdf/2602.18545

What happens in the background after the \code{quickCheck} call? The simplest form of a
property-based testing runner has two stages. First, a \generator{} produces values and
a \checker{} tests them, yielding one of three results: \passed{}, indicating that the
predicate did not reveal a bug on this input; \discarded{}, indicating that the input
is invalid for the property; or \failed{}, indicating that the input violates the
property. Once a failing input is found, the runner starts a second stage in which the
bug is minimized through shrinking, a form of
delta-debugging~\cite{zeller2002simplifying}. A \shrinker{} produces candidate smaller
values, for example by reducing keys in a tree or removing part of a computation. If a
candidate reproduces the failure, shrinking continues recursively from that candidate
until no smaller failing candidate remains, at which point the minimal value is
\printer{printed}. An abstract form of this process is depicted in
Figure~\ref{fig:quickcheck-property-runner-intro}:

\begin{figure}[h]
    \centering
    % TikZ styles for loop diagrams
% Shared styles for all interaction loop figures

% Component box base style
\tikzset{
  component/.style={
    minimum size=0.9cm,
    draw,
    line width=1.5pt,
    font=\Large\sffamily\bfseries,
    inner sep=0pt
  },
  % Individual component styles using paper's colors
  generator box/.style={component, draw=generator, text=generator},
  checker box/.style={component, draw=checker, text=checker},
  shrinker box/.style={component, draw=shrinker, text=shrinker},
  printer box/.style={component, draw=printer, text=printer, line width=1pt},
  mutator box/.style={component, draw=mutator, text=mutator},
  feedback box/.style={component, draw=feedback, text=feedback},
  combfeedback box/.style={component, draw=feedback, text=feedback, minimum size=0.7cm},
  size box/.style={component, draw=shrinker, text=shrinker, fill=shrinker!15, minimum size=0.8cm, font=\normalsize\bfseries},
  % Loop region style
  loop box/.style={
    draw=gray!70,
    dashed,
    line width=1pt,
    rounded corners=6pt,
    inner sep=8pt
  },
  loop region/.style={
    draw=gray!70,
    dashed,
    line width=1pt,
    rounded corners=6pt,
    fill=gray!8,
    inner sep=8pt
  },
  % Seed pool style
  seed pool/.style={
    draw=seedpool,
    dashed,
    line width=1pt,
    rounded corners=4pt,
    inner sep=5pt
  },
  % Arrow style
  loop arrow/.style={
    ->,
    >={Stealth[length=2.5mm]},
    line width=0.7pt
  },
  % Label style for arrows
  arrow label/.style={
    font=\tiny\sffamily,
    inner sep=1pt
  },
  % Loop title style
  loop title/.style={
    font=\scriptsize\sffamily,
    text=gray!70
  }
}

% Helper command for creating a standard G-C pair with bidirectional arrows
\newcommand{\gcpair}[3]{%
  % #1 = name prefix, #2 = x position, #3 = y position
  \node[generator box] (#1-G) at (#2, #3) {G};
  \node[checker box] (#1-C) at (#2+1.5, #3) {C};
  \draw[loop arrow] (#1-G.north) to[bend left=20] node[arrow label, above] {\texttt{check e}} (#1-C.north);
  \draw[loop arrow] (#1-C.south) to[bend left=20] node[arrow label, below] {\texttt{gen}} (#1-G.south);
}

% Helper command for shrinking loop (S-C pair)
\newcommand{\scpair}[3]{%
  % #1 = name prefix, #2 = x position, #3 = y position
  \node[shrinker box] (#1-S) at (#2, #3) {S};
  \node[checker box] (#1-C) at (#2+1.5, #3) {C};
  \draw[loop arrow] (#1-S.north) to[bend left=20] node[arrow label, above] {\texttt{check e'}} (#1-C.north);
  \draw[loop arrow] (#1-C.south) to[bend left=20] node[arrow label, below] {\texttt{shrink e | e'}} (#1-S.south);
}

    \resizebox{0.95\linewidth}{!}{% QuickChick interaction loop diagram
\begin{tikzpicture}[node distance=0.8cm]
  % Generation Loop
  \node[generator box] (G) {G};
  \node[checker box, right=1.5cm of G] (C1) {C};

  % Arrows in generation loop
  \draw[loop arrow] ([yshift=3pt]G.east) to[] node[arrow label, above] {\texttt{check e}} ([yshift=3pt]C1.west);
  \draw[loop arrow] ([yshift=-3pt]C1.west) to[] node[arrow label, below] {\texttt{gen}} ([yshift=-3pt]G.east);

  % Generation loop region
  \begin{scope}[on background layer]
    \node[loop region, fit=(G)(C1), inner ysep=16pt, inner xsep=6pt, label={[loop title]above:Generation Loop}] (genloop) {};
  \end{scope}

  % Shrinking Loop
  \node[shrinker box, right=1.5cm of C1] (S) {S};
  \node[checker box, right=1.5cm of S] (C2) {C};

  % Arrows in shrinking loop
  \draw[loop arrow] ([yshift=3pt]S.east) to[] node[arrow label, above] {\texttt{check e'}} ([yshift=3pt]C2.west);
  \draw[loop arrow] ([yshift=-3pt]C2.west) to[] node[arrow label, below] {\texttt{shrink e | e'}} ([yshift=-3pt]S.east);

  % Shrinking loop region
  \begin{scope}[on background layer]
    \node[loop region, fit=(S)(C2), inner ysep=16pt, inner xsep=6pt, label={[loop title]above:Shrinking Loop}] (shrinkloop) {};
  \end{scope}

  % Arrow from generation to shrinking
  \draw[loop arrow] (C1.east) -- node[arrow label, above] {\texttt{shrink e}} (S.west);

  % Printer
  \node[printer box, right=1.5cm of C2] (P) {P};

  % Arrow to printer
  \draw[loop arrow] (C2.east) -- node[arrow label, above] {\texttt{print e/e'}} (P.west);
\end{tikzpicture}
}
    \caption{An abstract representation of QuickCheck's property runner.}
    \label{fig:quickcheck-property-runner-intro}
\end{figure}

The implementations presented in the PBT frameworks table
(Table~\ref{tab:pbt-frameworks}) all implement different versions of this same core
algorithm, but vary the internal parameters of generators and shrinkers, the APIs for
expressing properties, and the underlying source of randomness. Hypothesis~\cite{Hypothesis},
Zest~\cite{Zest}, and Crowbar~\cite{Crowbar} parse randomness~\cite{goldsteinParsingRandomness2022}
into structures in their respective programming languages. Many other frameworks use
QuickCheck-style RNG splitting~\cite{ClaessenP13}, starting from an initial random seed
and splitting it at every decision point.

Many of these decisions are defaults that can be overridden when users need custom
configurations for their domains. Others \emph{leak} outside the abstraction and may
require significant reengineering of library internals to accommodate alternative
behaviors. For instance, the decision to use randomness parsing versus RNG splitting
affects one of the columns in the framework table
(Table~\ref{tab:pbt-frameworks}), \emph{integrated shrinking}. In integrated
shrinking~\cite{HypothesisShrinking}, in contrast to the classic structural shrinking
approach we mentioned, the source of randomness that is parsed to the generated
structure is shrunk instead of the structure itself. As such, a library using RNG
splitting cannot accommodate integrated shrinking without reengineering of its
internals.

This thesis argues that effective PBT frameworks require consolidation along two
axes: programmable interfaces that make techniques easier to express, combine,
and adopt; and shared evaluation infrastructure that grounds design choices in
bug-finding capability and quantitative measurement.

Chapter~\ref{chap:ppbt} on Programmable Property-Based
Testing~\cite{keles2026programmablepropertybasedtesting} focuses on this exact problem
and its many variants that cause fragmentation of the PBT landscape.
Table~\ref{tab:pbt-frameworks} presents 50 frameworks in 31 languages, with up to 5
libraries in OCaml. We trace a precise design decision, the design of the property
language in QuickCheck, as the predominant source of inflexibility of a PBT framework.
%
The shallow embedded domain-specific language forces the implementation of the property
runner to reside in the ``library-space'' as opposed to ``user-space'', leading to
significant ecosystem fragmentation.
%
We present an alternative property-language, a mixed-embedding, that enables flexible
user-space property runners, allowing us to combine many features and advances
dispersed throughout the ecosystem under a single library that users can keep on
extending.

Programmable PBT addresses a central challenge: consolidating algorithmic advances
within a single framework without reengineering internals, forking an existing library,
or writing a new library from scratch. It lowers the cost of experimentation and helps
new techniques diffuse through the ecosystem. However, fragmentation is not only a
matter of competing frameworks within a language; it is also a problem of information
dissemination. How can framework authors adopt a novel approach without concrete
empirical evidence about its effects in their context? A brief look at PBT papers
published in the last decade reveals a core issue: quantitative evaluation of PBT is
dispersed across bespoke, purpose-specific benchmarks built for individual papers and
rarely maintained afterward.

In Chapter~\ref{chap:etna}, I introduce ETNA~\cite{ETNA}, the first large-scale
evaluation platform for PBT. We designed ETNA to be easily extensible in terms of
adding new testing strategies for evaluation, new workloads as novel benchmarks, and
new languages with entirely separate ecosystems from the existing ones. The purpose of
ETNA is to consolidate knowledge through a common interface for quantitative
evaluation: shared benchmarks across languages for comparing cross-language
performance trade-offs, sampling-based evaluation modes that abstract away framework
specifics and focus on generated-test quality, and reproducible experiments for
sharing, verifying, and replicating existing results. We have
used ETNA to carry out experiments evaluating different generators, shrinkers, and
performance differences across frameworks and languages as we explain in
\S~\ref{evaluating-generators} and \S~\ref{evaluating-shrinking}. Similar to
Programmable PBT, ETNA is a step towards accelerating scientific development of PBT.

Much of the contemporary discussion of features, generalization, performance,
interface, and usage of property-based testing is based on general purpose frameworks
we have outlined in Table~\ref{tab:pbt-frameworks}. However, the world is full of
oddities, special-purpose tools, and requirements that push outside the boxes we draw.
General purpose tools grow as they redraw the lines of their approach to accommodate
such extensions. In a thesis focusing on consolidation, it is important to examine
what we cannot, at least for the moment, consolidate. In
Chapter~\ref{chap:dirt}, I present and discuss a domain-specific property-based testing
paradigm for databases, ``DIRT: Database-Integrated Random
Testing''~\cite{dirt-dbtest26}. Modern databases have enjoyed the results of fuzzing,
at times with property-based testing flavored approaches, in strengthening their
reliability and correctness. However, a major challenge that is yet to be solved is the
ability to evolve a property-based testing suite alongside the database, which we
attempted to solve. The chapter discusses various other domain-specific tools and
techniques as well as their potential inclusion to general purpose PBT frameworks.

The thesis concludes by returning to the central question running through these
chapters: how can PBT move from a collection of isolated framework traditions toward a
shared engineering and research discipline? The answer I develop is necessarily
partial. Flexible property languages, common evaluation infrastructure, and
domain-specific testing case studies do not settle every design question, but they make
those questions explicit and measurable.
